% !TEX root = immunology.tex

\section{Complement System}
  \label{sec:complement}
  
  \subsection{MAC Complex Attack}

    \begin{core}{Complement System}

      The \textbf{complement system} is an enzymatic cascade of more than 30 soluble plasma proteins and membrane receptors that function as a core branch of \textbf{innate immunity} while also amplifying the effector functions of adaptive antibodies.

      Most circulating complement proteins are synthesized predominantly by \textbf{hepatocytes} (\textit{liver cells}) and released as \textbf{inactive} precursors, although several immune and stromal cell populations also produce complement components locally.

    \end{core}

    \begin{core}[label=core:mac_complex]{MAC Complex Attack}

      The \textbf{membrane attack complex (MAC)} is the shared \textbf{terminal} effector of complement activation. The \textbf{classical} (\textit{\autoref{core:complement_classical}}), \textbf{lectin} (\textit{\autoref{core:complement_lectin}}), and \textbf{alternative} (\textit{\autoref{core:complement_alternative}}) pathways all converge on the generation of a \textbf{C5 convertase}, whose cleavage of C5 initiates assembly of the terminal complement components.

      \pgfig{
        [
          >=Stealth,
          every node/.style={font=\sffamily},
          component/.style={
            draw=black,
            fill=boxblue,
            rounded corners=7pt,
            minimum width=0.74cm,
            minimum height=1.55cm,
            inner sep=2pt,
            align=center
          },
          smallcomponent/.style={
            component,
            minimum width=0.68cm,
            minimum height=0.82cm
          },
          membrane/.style={draw=black, fill=boxblue},
          process/.style={->, thick},
          note/.style={font=\sffamily\small, align=left}
        ]

        % ===== Extracellular-space label =====
        \node[font=\sffamily\bfseries\large, anchor=west] at (0.35,10.85) {ECM};

        % ===== C5 cleavage =====
        \node[component] (c5) at (2.25,9.05) {C5};
        \node[smallcomponent] (c5a) at (0.95,6.95) {C5a};
        \node[component] (c5b) at (3.45,5.85) {C5b};

        \draw[process] (c5.south west) to[bend left=28] (c5a.east);
        \draw[process] (c5.south east) to[bend right=8] (c5b.north west);
        \node[font=\sffamily\small, rotate=-73] at (2.62,7.48) {Cleavage};

        \node[note, font=\sffamily\small, text width=1.65cm, anchor=north west] at (1,6.5)
          {(A)};

        \node[note, font=\sffamily\small, text width=1.65cm, anchor=north west] at (3.95,5.9)
          {(B)};

        % ===== Recruitment of C6 =====
        \node[component] (c6free) at (4.25,9.1) {C6};
        \node[component] (c5b6a) at (5.95,7.75) {C5b};
        \node[component] (c5b6b) at (6.55,8.02) {C6};

        \draw[process] (c5b.north east) -- (5.45,7.25);
        \draw[process] (c6free.south) -- (5.68,7.35);

        % ===== Recruitment of C7 and membrane attachment =====
        \node[component] (c7free) at (8.15,10.2) {C7};
        \draw[process] (c5b6b.east) -- (8.25,8.15);
        \draw[process] (c7free.south) -- (8.25,8.15);

        \draw[thick] (8.65,7.05) rectangle (11.65,9.95);
        \draw[membrane] (8.7,7.38) rectangle (11.6,7.82);
        \node at (10.15,7.6) {Cell Membrane};
        \node[component] at (9.45,8.52) {C5b};
        \node[component] at (10.02,8.72) {C6};
        \node[component] at (10.60,8.55) {C7};
        \node[note, text width=2.55cm, anchor=south west] at (9.05,10.05)
          {(C)};

        % ===== Recruitment of C8 and membrane insertion =====
        \node[component] (c8free) at (13.95,10.25) {C8};
        \coordinate (c8join) at (13.0,7.15);
        \draw[process] (11.78,8.55) to[bend left=18] (c8join);
        \draw[process] (c8free.south west) -- (c8join);

        \draw[thick] (12.15,3.75) rectangle (15.85,7.05);
        \draw[membrane] (12.2,4.55) rectangle (15.8,5.02);
        \node at (14.0,4.78) {Cell Membrane};
        \node[component] at (13.0,5.72) {C5b};
        \node[component] at (13.55,5.94) {C6};
        \node[component] at (14.10,5.78) {C7};
        \node[component, minimum height=2.05cm] at (14.68,5.35) {C8};
        \node[note, text width=2.9cm, anchor=north west] at (12.7,3.48)
          {(D)};

        % ===== C9 recruitment =====
        \node[component, minimum height=1.8cm] (c9free) at (10.55,3.25) {C9};
        \draw[process] (c9free.west) to[bend right=32] (8.78,3);
        \draw[process] (13.1,3.72) to[bend left=30] (8.15,2.48);
        \node[note, anchor=east] at (9.95,3.6) {C9 polymerization};

        % ===== Mature membrane attack complex =====
        \draw[membrane] (0.25,0.95) rectangle (15.65,1.42);
        \node[anchor=east] at (15.5,1.18) {Cell Membrane};

        \node[component] at (5.55,2.12) {C5b};
        \node[component] at (6.12,2.32) {C6};
        \node[component] at (6.69,2.18) {C7};
        \node[component, minimum height=2.05cm] at (7.25,1.73) {C8};

        \node[component, minimum height=2.55cm, inner sep=0pt] at (7.70,1.42) {\small C9};
        \node[component, minimum height=2.72cm, inner sep=0pt] at (8.16,1.42) {\small C9};
        \node[component, minimum height=2.85cm, inner sep=0pt] at (8.62,1.42) {\small C9};
        \node[component, minimum height=2.72cm, inner sep=0pt] at (9.08,1.42) {\small C9};
        \node[component, minimum height=2.55cm, inner sep=0pt] at (9.54,1.42) {\small C9};

        \draw[-{Stealth[length=9pt,width=10pt]}, line width=3.2pt]
          (8.62,3.25) -- (8.62,0.15);
        \node[font=\sffamily\bfseries\large, anchor=west] at (9.95,1.72) {MAC Pore};
        \node[anchor=north] at (8.62,0.08) {Water and ions};
        \node[font=\sffamily\bfseries\large, anchor=east] at (15.4,0.1)
          {Inside of Target Cell};
      }{Membrane attack complex (MAC); \textit{all different complement pathways converge to this pathway}; (A) \textit{C5a acts as a strong \textbf{anaphylatoxin} (\autoref{core:side_pathways})}; (B) \textit{C5b is unstable when left unbound}; (C) \textit{C7 recruitment exposes a hydrophobic site that anchors the complex to the membrane}; (D) \textit{C8 recruitment drives partial insertion of the complex into the lipid bilayer}}{fig:mac}{0.9}
      \vspace{-0.4cm}

      As shown in \autoref{fig:mac}, the completed MAC achieves \textbf{membrane lysis} by forming a transmembrane pore composed of polymerized C9 subunits; structural studies generally describe a pore containing approximately 18 C9 subunits. This pore permits the \textbf{uncontrolled movement of water and ions across the membrane}, disrupting ionic homeostasis and potentially inducing osmotic rupture. Several types of susceptible targets can be damaged in this way, including the following.

      \begin{itemize}

        \item \textbf{Enveloped viruses}, whose host-derived lipid membranes may be disrupted by terminal complement deposition

        \item \textbf{Gram-negative bacteria} (\textit{FIGURE REFEERNCE FROM UNIT 1}), whose outer membrane is accessible to MAC insertion, in contrast to the thick peptidoglycan wall that generally protects \textit{Gram-positive bacteria} from direct MAC-mediated lysis

        \item Susceptible \textbf{eukaryotic parasites}

        \item \textbf{Erythrocytes} (\textit{red blood cells}) under pathological conditions in which complement regulation fails or inappropriate activation occurs

        \item Antibody-targeted \textbf{cancer cells}, provided that complement activation overcomes the cells' membrane-bound complement regulators.

      \end{itemize}

    \end{core}

  \subsection{Complement Pathways}

    \begin{core}[label=core:complement_classical]{Classical Pathway}

      The \textbf{classical pathway} is the antibody-dependent route of complement activation, triggered when the C1 complex engages antigen-bound immunoglobulin.

      \myfig{0.5}{immunology_figures/c1_complex}{C1 complex attached to the Fc stem of an IgG bound to an antigen}{fig:c1_complex}

      As shown in \autoref{fig:c1_complex}, the inactive C1 complex is composed of one hexameric C1q recognition molecule bound to two C1r and two C1s serine proteases. Each of C1q's six arms terminates in a globular head assembled from C1qA, C1qB, and C1qC chains (\textit{not shown}). C1q recognizes sites within the Fc region of antigen-bound immunoglobulin, and the local Fc conformation and associated glycans influence the accessibility and efficiency of this interaction. \textbf{Engagement of multiple Fc regions} by C1q induces a conformational change that activates C1r, which in turn activates C1s, initiating proteolytic cleavage of C4 and C2.

      \pgfig{
        [
          >=Stealth,
          every node/.style={font=\sffamily},
          component/.style={
            draw=black,
            fill=boxblue,
            rounded corners=7pt,
            minimum width=0.72cm,
            minimum height=1.45cm,
            inner sep=2pt,
            align=center
          },
          fragment/.style={
            component,
            rounded corners=5pt,
            minimum width=0.65cm,
            minimum height=0.82cm
          },
          process/.style={->, thick},
          note/.style={font=\sffamily\small, align=left}
        ]

        % ===== Activation of the C1 serine proteases =====
        \node[font=\sffamily\bfseries\large] (c1r) at (14.0,10.6) {C1r};
        \node[font=\sffamily\bfseries\large] (c1s) at (11.75,10.6) {C1s};
        \draw[process] (c1r.west) -- (c1s.east);
        \node[note, font=\sffamily\small, text width=2.7cm, anchor=north west]
          at (12.25,10.15) {(A)};

        % ===== C2 and C4 cleavage =====
        \node[component] (c2) at (5.0,8.9) {C2};
        \node[component] (c4) at (9.7,8.9) {C4};
        \draw[process] (c1s.south west) -- (c2.north east)
          node[midway, above, sloped, font=\sffamily\small] {cleaves};
        \draw[process] (c1s.south) -- (c4.north east)
          node[midway, below, sloped, font=\sffamily\small] {cleaves};

        \node[fragment] (c2b) at (4.55,6.85) {C2b};
        \node[fragment] (c2a) at (6.35,6.85) {C2a};
        \draw[process] (c2.south west) -- (c2b.north);
        \draw[process] (c2.south east) -- (c2a.north);

        \node[fragment] (c4b) at (9.1,6.85) {C4b};
        \node[fragment] (c4a) at (11.0,6.85) {C4a};
        \draw[process] (c4.south west) -- (c4b.north);
        \draw[process] (c4.south east) -- (c4a.north);
        \node[note, text width=2.5cm, anchor=west] at (11.55,6.85)
          {(B)};

        % ===== Classical C3 convertase =====
        \node[component] (c4bconv) at (7.98,4.95) {C4b};
        \node[fragment, minimum height=1.15cm] (c2aconv) at (7.25,4.95) {C2a};
        \draw[process] (c4b.south west) -- (c4bconv.north east);
        \draw[process] (c2a.south east) -- (c2aconv.north west);
        \node[font=\sffamily\bfseries, anchor=north] at (10.62,5.18) {C4b2a (\textit{C3 convertase})};

        % ===== C3 cleavage and effector fragments =====
        \node[component] (c3) at (2.3,4.95) {C3};
        \draw[process] (c2aconv.west) -- (c3.east)
          node[midway, above, font=\sffamily\small] {cleaves};
        \node[fragment] (c3a) at (1.85,2.85) {C3a};
        \node[fragment] (c3b) at (3.55,2.85) {C3b};
        \draw[process] (c3.south west) -- (c3a.north);
        \draw[process] (c3.south east) -- (c3b.north);
        \node[note] at (1.45,1.8) {(C)};
        \node[note] at (3.5,1.8) {(D)};

        % ===== Classical C5 convertase =====
        \node[component] (c4bc5) at (8.62,1.75) {C4b};
        \node[fragment, minimum height=1.15cm] (c2ac5) at (7.88,1.75) {C2a};
        \node[component] (c3bc5) at (7.15,1.75) {C3b};
        \draw[process] (7.55,4.2) -- (8,2.5);
        \draw[process] (c3b.east) -- (c3bc5.north west);
        \node[font=\sffamily\bfseries, anchor=north] at (7.88,0.98) {C4b2a3b (\textit{C5 convertase})};

        % ===== Entry into the terminal pathway =====
        \node[font=\sffamily\bfseries\large, align=center] (mac) at (13.1,1.8)
          {MAC Activation\\Pathway};
        \draw[process] (c4bc5.east) -- (mac.west)
          node[midway, above, font=\sffamily\small\itshape] {cleaves C5};
      }{Classical complement pathway activation cascade; (A) \textit{activated C1r cleaves C1s, producing an active protease}; (B) \textit{C4a is a weak \textbf{anaphylatoxin} (\autoref{core:side_pathways})}; (C) C3a goes onto act as an \textbf{anaphylatoxin}; (D) \textit{C3b also functions as an \textbf{opsonin} (\autoref{core:side_pathways}), coating targets for phagocytic recognition}}{fig:classical_pathway}{0.9}
      \vspace{-0.4cm}

      It's important to note that complement-activating antibody isotypes differ substantially in their ability to recruit C1q.

      \begin{itemize}

        \item \textbf{IgM} is a \textbf{highly efficient} activator. Soluble pentameric IgM generally adopts a planar conformation, whereas when bound to antigen, the pentamer takes up a \textbf{staple-like conformation} (\textit{FIG REFERENCE}) that exposes C1q-binding sites in its Fc regions.

        \myfig{0.9}{immunology_figures/igm_staple}{Antigen binding induces a conformational shift in pentameric IgM, from a planar resting state (\textit{left}) to a staple-like conformation (\textit{right}) that exposes C1q-binding sites in the Fc stems}{fig:igm_staple}

        \item Because \textbf{IgG} is monomeric, efficient C1 activation normally requires several antigen-bound IgG molecules to cluster their Fc regions and provide multivalent engagement of C1q. This requirement limits inappropriate activation by freely circulating IgG.

        Among human IgG subclasses, classical pathway activation potency varies considerably: IgG3, owing to its comparatively long hinge region, is the most efficient activator, followed by IgG1, while IgG2 is relatively weak and IgG4 does not appreciably activate the classical complement pathway.

        \item \textbf{IgA, IgE, and IgD} do not ordinarily initiate the classical complement pathway.

      \end{itemize}

    \end{core}

    \begin{core}[label=core:complement_lectin]{Lectin Pathway}

      The \textbf{lectin pathway} parallels the classical pathway downstream of its initial recognition event. Instead of using C1q to recognize antigen-bound antibody, it begins when soluble \textbf{PRR} molecules, namely \textbf{mannose-binding lectin (MBL)}, bind characteristic carbohydrate structures on a microbial surface and activate associated serine proteases called \textbf{MASP}.

      \myfig{0.5}{immunology_figures/mbl_masp}{Structure of tetrameric MBL bound to a MASP dimer}{fig:mbl_masp}

      In comparison to the \textbf{classical pathway} (\textit{\autoref{core:complement_classical}}), MBL resembles C1q in that it occurs in clusters. Additionally MASP molecules are anlogous to C1r and C1s

      \begin{table}[H]
        \centering
        \begin{tabular}{>{\raggedright\arraybackslash}p{6cm}>{\raggedright\arraybackslash}p{6cm}}
          \hline
          \textbf{Lectin Pathway} & \textbf{Classical Pathway} \\
          \hline
          Initiated by mannose and other carbohydrates & Initiated by antibodies, using the domain with the oligosaccharide \\
          Changes conformation of MBL & Changes conformation of C1q \\
          Activates MASP & Activates C1r/C1s \\
          \hline
        \end{tabular}
        \caption{Comparison of the lectin and classical complement pathways}
        \label{tab:lectin_vs_classical}
      \end{table}

      Following activation of the lectin-pathway proteases, cleavage of C4 and C2 produces the same C3 convertase, C4b2a, that functions in the classical pathway. The two pathways are therefore identical from C3-convertase formation onward.

    \end{core}

    \begin{core}[label=core:complement_alternative]{Alternative Pathway}

      The \textbf{alternative pathway} is an evolutionarily ancient route of complement activation that relies on innate pattern recognition of antigens present on pathogenic surfaces, including teichoic acid (\textit{bacterial}), zymosan (\textit{fungal}), certain tumor-associated compounds, and trypanosome (\textit{unicellular eukaryote}) surface markers. Like the classical (\textit{\autoref{core:complement_classical}}) and lectin (\textit{\autoref{core:complement_lectin}}) pathways, the alternative pathway ultimately converges on the formation of a C5 convertase, which cleaves C5 and initiates assembly of the MAC, as outlined in \autoref{core:mac_complex}.

      \pgfig{
        [
          >=Stealth,
          every node/.style={font=\sffamily},
          component/.style={
            draw=black,
            fill=boxblue,
            rounded corners=7pt,
            minimum width=0.76cm,
            minimum height=1.45cm,
            inner sep=2pt,
            align=center
          },
          fragment/.style={
            component,
            rounded corners=5pt,
            minimum width=0.68cm,
            minimum height=0.82cm
          },
          factor/.style={draw=black, fill=boxblue, minimum width=1.75cm, minimum height=0.72cm},
          process/.style={->, thick},
          note/.style={font=\sffamily\small, align=left}
        ]

        % ===== Spontaneous C3 tickover =====
        \node[component] (c3tick) at (0.95,8.55) {C3};
        \node[font=\sffamily\bfseries\large] (water) at (3.25,10.2) {\ce{H2O}};
        \node[note, text width=2.0cm, anchor=north west] at (0.3,10.6)
          {Spontaneous hydrolysis};
        \draw[process] (water.west) to[bend right=18] (c3tick.north east);

        \node[component, align=center] (c3water) at (3.6,8.35) {C3\\\ce{H2O}};
        \draw[process] (c3tick.east) -- (c3water.west);
        \node[font=\sffamily\small\bfseries, align=center, anchor=north]
          at (3.6,7.55) {Hydrolyzed C3\\C3(\ce{H2O})};

        % ===== Factor B binding and Factor D cleavage =====
        \node[factor] (factorb) at (5.45,9.95) {Factor B};
        \node[component, align=center] (c3waterb) at (7.25,8.35) {C3\\\ce{H2O}};
        \node[factor, minimum width=1.3cm, minimum height=1.0cm] (boundb) at (8.35,8.35) {Factor B};
        \draw[process] (c3water.east) -- (c3waterb.west);
        \draw[process] (factorb.south east) -- (c3waterb.north west);

        \node[factor] (factord) at (11.25,9.75) {Factor D};
        \draw[process] (factord.south west) -- (boundb.north east)
          node[midway, right, font=\sffamily\small, yshift=-0.15cm] {cleaves};
        \node[factor, minimum width=1.25cm] (ba) at (11.85,7.55) {Factor Ba};
        \draw[process] (boundb.east) -- (ba.west);
        \node[note, font=\sffamily\small\itshape, anchor=west] at (12.65,7.55)
          {(A)};

        % ===== Fluid-phase C3 convertase =====
          \node[component, align=center] (c3waterconv) at (7.55,5.85) {C3\\\ce{H2O}};
        \node[factor, minimum width=1.15cm, minimum height=0.7cm] (bbfluid) at (8.65,5.85) {Bb};
        \draw[process] (c3waterb.south) -- (c3waterconv.north);
        \node[font=\sffamily\bfseries, anchor=north] at (8.05,5.08) {C3(\ce{H2O})Bb};
        \node[font=\sffamily\small\bfseries, anchor=north] at (8.05,4.62)
          {(fluid C3 convertase)};

        % ===== Amplification through C3 cleavage =====
        \node[component] (c3substrate) at (3.45,5.8) {C3};
        \draw[process] (c3waterconv.west) -- (c3substrate.east)
          node[midway, above, font=\sffamily\small] {cleaves};
        \node[fragment] (c3aalt) at (2.25,3.85) {C3a};
        \node[fragment] (c3balt) at (4.55,3.85) {C3b};
        \draw[process] (c3substrate.south west) -- (c3aalt.north);
        \draw[process] (c3substrate.south east) -- (c3balt.north);
        \node[note] at (1.8,2.85)
          {(B)};
        \node[note] at (4.55,2.85)
          {(C)};

        % ===== Surface C3 convertase and properdin stabilization =====
        \node[component] (c3bsurface) at (7.55,2.55) {C3b};
        \node[factor, minimum width=1.05cm, minimum height=0.72cm] (bbsurface) at (8.55,2.55) {Bb};
        \draw[process] (c3balt.east) -- (c3bsurface.west);
        \node[note, anchor=west] at (6,3.42) {(D)};
        \node[font=\sffamily\bfseries, anchor=north] at (8.05,1.78) {C3bBb};
        \node[font=\sffamily\small\bfseries, anchor=north] at (8.05,1.32)
          {(surface C3 convertase)};

        \node[factor, align=center] (properdin) at (12.65,5.25) {Properdin\\(Factor F)};
        \draw[process, dashed] (properdin.west) to[bend right=12] (bbsurface.north east);
        \node[note, font=\sffamily\small\itshape, text width=3.4cm, anchor=north west]
          at (10.4,4.45) {(E)};

        % ===== Alternative C5 convertase =====
        \node[fragment, minimum width=0.85cm] (c3bextra) at (10.65,2.45) {C3b};
        \node[note, anchor=west] at (11.15,2.45) {Additional C3b};

        \node[component] (c3bc5a) at (9.55,-0.65) {C3b};
        \node[factor, minimum width=1.0cm, minimum height=0.72cm] (bbc5) at (10.55,-0.65) {Bb};
        \node[component] (c3bc5b) at (11.55,-0.65) {C3b};
        \draw[process] (bbsurface.east) to[bend left=28] (c3bc5a.north east);
        \draw[process] (c3bextra.south) -- (c3bc5b.north);
        \node[font=\sffamily\small\bfseries, anchor=north] at (10.55,-1.43)
          {C3bBbC3b (C5 convertase)};

        % ===== Entry into the terminal pathway =====
        \node[font=\sffamily\bfseries\large, align=center] (macalt) at (15.6,-0.55)
          {MAC Activation\\Pathway};
        \draw[process] (c3bc5b.east) -- (macalt.west);
        \node[font=\sffamily\small\itshape] at (12.8,-0.5) {cleaves C5};
      }{Alternative complement pathway activation cascade; (A) \textit{factor Ba diffuses away}; (B) \textit{C3a goes onto act as an \textbf{anaphylatoxin} (\autoref{core:side_pathways})}; (C) \textit{C3b also serves as an \textbf{opsonin} (\autoref{core:side_pathways}), coating activating surfaces}; (D) \textit{factor B binding and factor D cleavage of factor B into Bb}; (E) \textit{properdin---sometimes called factor F--- binds and stabilizes C3bBb on the activating surface}}{fig:alternative_pathway}{0.9}

    \end{core}

    \begin{core}[label=core:side_pathways]{Side Signal Pathways \& Other}

      Complement activation produces several effects in addition to MAC-mediated lysis. Cleavage products generated throughout the cascade promote inflammation, coat targets for immune recognition, neutralize pathogens, and facilitate the removal of immune complexes.

      An \textbf{anaphylatoxin} is a small peptide fragment released as a byproduct of proteolytic cleavage during the complement cascade; upon binding specific receptors on immune and vascular cells, these fragments trigger localized inflammatory responses. C3a, C4a, and C5a are all anaphylatoxins, arising respectively from cleavage of C3, C4, and C5 during the unique complement and converged MAC pathways.

      \begin{itemize}

        \item \textbf{C3a}: C3a stimulates degranulation of \textbf{mast cells} and \textbf{basophils}, promoting the release of \textbf{histamine} and other inflammatory mediators that increase vascular permeability.
          
        \item \textbf{C4a}: C4a exhibits \textbf{comparatively weak anaphylatoxin activity} relative to C3a and C5a, but nonetheless contributes to local inflammatory signaling following its release during C4 cleavage.

        \item \textbf{C5a}: Among the anaphylatoxins, C5a is the most potent, functioning as a chemoattractant that recruits and activates \textbf{neutrophils}, \textbf{monocytes}, and other \textbf{myeloid cells} at sites of complement activation.

      \end{itemize}

      In addition to their inflammatory functions, complement fragments deposited on target surfaces mark pathogens and immune complexes for recognition, neutralization, and clearance by phagocytic mechanisms.

      \begin{enumerate}

        \item \textbf{Opsonization}: Covalently deposited complement fragments, particularly C3b and iC3b, coat microbial surfaces and enhance their recognition and phagocytosis through complement receptors.

        \item \textbf{Viral Neutralization}: Complement deposition can sterically hinder viral attachment or entry, promote virion aggregation, and facilitate clearance by phagocytes.

        \item \textbf{Immune-Complex Clearance}: Complement-tagged antigen--antibody complexes bind complement receptor 1 on erythrocytes and are transported through the circulation to the liver and spleen, where phagocytes remove and degrade them.

      \end{enumerate}

    \end{core}

\section{Major Histocompatibility Complex}

  \subsection{MHC Structure \& Genetics}

    \begin{core}[label=core:mhc_characteristics]{MHC Characteristics}

      The \textbf{Major Histocompatibility Complex (MHC)} is a set of closely linked genes that encode cell-surface glycoproteins essential for the adaptive immune system. Its primary function is to \textbf{bind antigens} derived from pathogens or self-proteins and \textbf{display them on the cell surface} for recognition by \textbf{T lymphocytes} via \textbf{T-cell receptors} (\textit{\autoref{sec:tcr}}).

      \pgfig{
        [
          >=Stealth,
          mhcdomain/.style={ellipse, draw=none, minimum width=2.15cm, minimum height=2.7cm,
            inner sep=0pt, font=\bfseries\large, align=center},
          annotation/.style={font=\scriptsize, align=center},
          heading/.style={font=\bfseries},
          pointer/.style={-{Stealth[length=6pt,width=6pt]}, line width=0.6pt},
          anchorline/.style={line width=5.5pt, line cap=round}
        ]

        \definecolor{mhcired}{RGB}{220,0,0}
        \definecolor{mhciired}{RGB}{226,99,102}
        \definecolor{mhciiblue}{RGB}{103,151,224}
        \definecolor{mhcblue}{RGB}{61,120,211}
        \definecolor{cdgreen}{RGB}{101,171,76}
        \definecolor{peptidemagenta}{RGB}{241,0,226}
        \definecolor{membraneblue}{RGB}{207,226,243}

        % ===== Compartments and cell membrane =====
          \node[font=\bfseries\Large, anchor=east] at (18.45,8) {ECM};
          \filldraw[fill=membraneblue, draw=black, line width=0.65pt]
            (0,1.05) rectangle (18.55,1.68);
          \node[font=\bfseries\large, anchor=east] at (18.35,1.36)
            {Cell Membrane};
          \node[font=\bfseries\Large, anchor=east] at (18.45,0.37)
            {Cytoplasm};

        % ===== MHC class I =====
          \node[font=\bfseries\large] at (3.35,8.5) {[MHC I]};

          \node[] at (5.2,4.08) {(B)};

          \node[] at (3.35,7.7) {(A)};

          \node[mhcdomain, fill=mhcired] (ia2) at (2.45,5.63) {α2};
          \node[mhcdomain, fill=mhcired] (ia1) at (4.23,5.63) {α1};
          \node[mhcdomain, fill=mhcired] (ia3) at (2.45,3.02) {α3};
          \node[mhcdomain, fill=mhcblue, minimum width=1.72cm, minimum height=2.38cm]
            (ibeta) at (4.27,3.02) {β};

          \draw[peptidemagenta, line width=4pt, line cap=round]
            (3.08,6.86) .. controls (3.16,7.30) and (3.53,7.30) .. (3.65,6.88);
          \draw[anchorline, mhcired] (2.45,1.76) -- (2.45,0.98);
          \fill[mhcired] (2.45,0.98) circle (0.18);
          \node[] at (2.45,0.3) {(C)};

        % ===== MHC class II =====
          \node[font=\bfseries\large] at (9.2,8.5) {[MHC II]};

          \node[] at (11.3,6.96) {(E)};

          \node[] at (9.2,7.7) {(D)};

          \node[mhcdomain, fill=mhciired] (iia1) at (8.08,5.63) {α1};
          \node[mhcdomain, fill=mhciired] (iia2) at (8.08,3.02) {α2};
          \node[mhcdomain, fill=mhciiblue] (iib1) at (10.28,5.63) {β1};
          \node[mhcdomain, fill=mhciiblue] (iib2) at (10.28,3.02) {β2};
          \draw[peptidemagenta, line width=4pt, line cap=round]
            (9.10,5.45) .. controls (9.28,5.70) and (9.27,6.31) .. (9.35,6.55);
          \draw[anchorline, mhciired] (8.08,1.76) -- (8.08,0.98);
          \draw[anchorline, mhciiblue] (10.28,1.76) -- (10.28,0.98);
          \fill[mhciired] (8.08,0.98) circle (0.18);
          \fill[mhciiblue] (10.28,0.98) circle (0.18);

        % ===== CD1 (MHC-I-like) molecule =====
          \node[font=\bfseries\large] at (14.6,8.5) {[CD1 MHC]};

          \node[] at (14.6,7.7) {(F)};

          \node[] at (16.6,7) {(G)};

          \node[mhcdomain, fill=cdgreen] (cda2) at (13.75,5.63) {α2};
          \node[mhcdomain, fill=cdgreen] (cda1) at (15.55,5.63) {α1};
          \node[mhcdomain, fill=cdgreen] (cda3) at (13.75,3.02) {α3};
          \node[mhcdomain, fill=mhcblue, minimum width=1.72cm, minimum height=2.38cm]
            (cdbeta) at (15.59,3.02) {β};
          \draw[peptidemagenta, line width=4pt, line cap=round]
            (14.72,6.12) .. controls (14.86,6.45) and (14.84,6.83) .. (14.93,7.25);
          \draw[anchorline, cdgreen] (13.75,1.76) -- (13.75,0.98);
          \fill[cdgreen] (13.75,0.98) circle (0.18);

      }{MHC I (\textit{left}), MHC II (\textit{middle}), and CD1 MHC (\textit{right}) glycoproteins. (A) \textit{closed cleft, bowed antigen-presentation surface (8–10 amino acids; displays peptide middle)}; (B) \textit{β peptide linked to α1 and membrane by weak bonds, required for membrane expression}; (C) \textit{α3 connects to transmembrane segment with cytoplasmic tail}; (D) \textit{open cleft, flat antigen-presentation surface (13–18 amino acids, only 13 fit the cleft)}; (E) \textit{α/β dimer joined by weak interactions}; (F) \textit{deep cleft, flat lipid-presentation surface}; (G) \textit{closely resembles MHC I; β subunit is the identical peptide}}{fig:mhc}{0.8}

      As shown in \autoref{fig:mhc}, three distinct types of MHC receptors exist, each characterized by a unique combination of function, expression pattern, subunit composition, and binding site specificity; this information is summarized in the following table.

      \begin{table}[H]
        \centering
        \footnotesize
        \begin{tabular}{>{\raggedright\arraybackslash}p{1.2cm}>{\raggedright\arraybackslash}p{5cm}>{\raggedright\arraybackslash}p{4cm}>{\raggedright\arraybackslash}p{2.7cm}}
          \hline
          \textbf{} & \textbf{MHC I (Identification)} & \textbf{MHC II (Presentation)} & \textbf{CD1 MHC (Non-Classical MHC)} \\
          \hline
          \textbf{Function} &
          Displays \textbf{endogenous} peptide/antigen (\textit{\autoref{core:endogenous_mhci}}) to T\textsubscript{C} cells as a means of immune surveillance, determining whether the T cell will induce killing. Via \textbf{cross-presentation} (\textit{\autoref{core:cross_presentation}}), APCs also display \textbf{exogenous} peptides/antigens to T\textsubscript{C} cells to activate them. &
          Displays \textbf{exogenous} peptide/antigen (\textit{\autoref{core:exogenous_mhcii}}) to T\textsubscript{H} cells. &
          Displays \textbf{exogenous} bacterial glycolipids (\textit{\autoref{core:exogenous_cd1mhc}}) to T\textsubscript{H} cells. \\
          \textbf{Expression} &
          Expressed on \textbf{most nucleated cells}, with the exception of RBCs, sperm, and nerve cells. Expression is \textbf{highest on lymphocytes} and comparatively low on fibroblasts, muscle cells, and liver cells. &
          Expressed on \textbf{APCs} and on \textbf{thymic epithelial cells}, where it mediates T cell selection; other cell types can be induced to express it in response to certain cytokines. &
          Expressed on \textbf{APCs}. \\
          \textbf{Subunits} &
          [1 peptide, 3 domains] + [β globulin] &
          [2 peptides, 4 domains] &
          [1 peptide, 3 domains] + [β globulin] \\
          \textbf{Binding Site} &
          Consists of a \textbf{closed cleft} between α1 and α2, with α-helical sides and a β-sheet floor. Amino acid \#9 of the antigen peptide (\textit{C-terminal, occasionally assisted by \#8 and \#7}) anchors into the cleft's hydrophobic pocket, while amino acid \#2 (\textit{N-terminal}) anchors at the opposite end; the intervening, displayed residues are what the TCR ultimately engages.[\autoref{fig:mhc_structure}] &
          Consists of an \textbf{open cleft} between α1 and β1, with α-helical sides and a β-sheet floor; peptides can engage MHC II at multiple points along the cleft. [\autoref{fig:mhc_structure}] &
          Consists of a deep cleft between α1 and α2. \\
          \textbf{Antigen Binding Site} &
          Antigen binding occurs in the RER (\textit{recall that MHC I displays endogenous peptides}). &
          Antigen binding occurs in the phagolysosome. &
          Antigen binding occurs in the phagolysosome. \\
          \textbf{Versions} &
          3 versions (A, B, C) &
          3 versions (DP, DQ, DR) &
          5 versions (A, B, C, D, E) \\
          \textbf{Chromosome} &
          Chromosome 6 &
          Chromosome 6 &
          Chromosome 1 \\
          \hline
        \end{tabular}
        \caption{Comparison of MHC I, MHC II, and CD1 MHC}
        \label{tab:mhc_comparison}
      \end{table}

      \myfig{0.7}{immunology_figures/mhc_structure}{Three-dimensional structures of MHC I (\textit{left}) and MHC II (\textit{right}); \textit{both peptide-binding clefts are formed by α-helical walls flanking a β-sheet floor}}{fig:mhc_structure}

    \end{core}

    \begin{core}{MHC Gene Sequences}

      The \textbf{human MHC gene} is encoded within the short arm of \textbf{chromosome 6}, where the linked set of MHC alleles inherited together on a single chromosome constitutes an \textbf{MHC haplotype}; each diploid individual inherits one such haplotype maternally and one paternally, with both sets of alleles expressed \textbf{codominantly} (\textit{\autoref{gen-core:single_gene_abnormal}}). The corresponding gene products are designated \textbf{human leukocyte antigens (HLAs)} and, while products generally adopt immunoglobulin-superfamily folds, the MHC region also encodes numerous genes whose products fall outside this structural family, most notably within the MHC class III gene region, though a smaller number are also found interspersed among the class I and class II loci.

      \myfig{0.7}{immunology_figures/hla_gene}{Organization of the HLA (MHC) region on chromosome 6}{fig:hla_gene}

      \begin{enumerate}

        \item \textbf{MHC I} (\textit{HLA-A, HLA-B, HLA-C})

          Classical MHC class I molecules are encoded by the \textbf{HLA-A, HLA-B, and HLA-C} loci, which give rise to allelic products differing in peptide-binding specificity; these \textbf{loci encode only the polymorphic α chain}, while the invariant \textbf{β} component is encoded separately by the \textit{B2M} gene on chromosome 15. The HLA-A, HLA-B, and HLA-C loci are exceptionally \textbf{polymorphic}---modern allele databases list many thousands of named variants---such that an individual heterozygous at all three loci may express as many as six distinct classical MHC class I allotypes.

          This extraordinary allelic polymorphism is concentrated particularly within the sequences encoding the \textbf{peptide-binding groove}, and the resulting diversity among individuals primarily reflects inherited germline polymorphism rather than the somatic DNA rearrangement responsible for generating diversity in the CDRs of immunoglobulins and T-cell receptors.

        \item \textbf{MHC III}

          Unlike MHC classes I and II, the MHC class III region \textbf{does not encode peptide-presenting molecules}; it is classified as part of the MHC solely on the basis of its physical location between the class I and class II regions on chromosome 6. Instead, this region contains genes encoding a range of immune mediators, including \textbf{complement components} (\textit{\autoref{sec:complement}}), \textbf{tumor necrosis factors} (\textit{\autoref{core:tnf}}), and heat-shock proteins.

        \item \textbf{MHC II} (\textit{HLA-DP, HLA-DQ, HLA-DR})

          The classical MHC class II region contains the \textbf{HLA-DP, HLA-DQ, and HLA-DR} loci, each encoding an \textbf{αβ heterodimer} in which both chains contribute to the peptide-binding groove.

          Codominant expression of maternal and paternal class II genes, together with permitted pairing between certain α and β chains, \textbf{allows an individual to display a diverse set of MHC class II molecules}.

          As with class I, classical MHC class II loci are \textbf{highly polymorphic}, with much of this variation concentrated in residues that shape the peptide-binding groove, and the resulting interindividual diversity likewise reflects inherited germline polymorphism rather than somatic DNA rearrangement.

        \item \textbf{Other}

          The class II region also contains genes encoding components of the MHC class I antigen-processing pathway, including \textit{TAP1}, \textit{TAP2}, \textit{PSMB8}, and \textit{PSMB9}, while nonclassical MHC class I molecules themselves perform more specialized functions in immune recognition, molecular transport, and antigen presentation.

      \end{enumerate}

    \end{core}
    
    \begin{core}{MHC Mendalian Genetics}

      Because the MHC comprises a cluster of tightly \textbf{linked} (\textit{\autoref{gen-core:gene_linkage}}) loci rather than a single gene, its pattern of inheritance follows the segregation of whole haplotypes and is shaped by the same principles of linkage, recombination, and inbreeding that govern Mendelian genetics more broadly.

      \begin{enumerate}

        \item \textbf{Non-Recombinant Inheritance} 

          Due to the MHC loci being tightly linked, their alleles are usually inherited together as a single haplotype, so that in the absence of recombination two full siblings have a 25\% probability of inheriting the same pair of parental MHC haplotypes.

          \myfig{0.45}{immunology_figures/mhc_nonrecomb}{Non-recombinant HLA gene inheritance; \textit{offspring receive a random MHC haplotype from each parent}}{fig:mhc_nonrecomb}

          \textit{More generally, the term \textbf{haplotype} denotes a set of linked alleles or sequence variants inherited together from one parent; its use is not restricted to the MHC.}

        \item \textbf{Recombinant Inheritance}

          Meiotic \textbf{crossing over} (\textit{\autoref{gen-core:crossing_over}}) within the MHC can generate a recombinant haplotype that combines portions of the parental class I and class II regions, thereby reducing the probability that siblings will match across the entire MHC; over evolutionary time, this same process contributes additional haplotypic diversity to the population.

          Conversely, \textbf{inbreeding} increases the probability that an individual will inherit identical or closely related MHC haplotypes; in a highly inbred population, the limited introduction of new haplotypes increases MHC homozygosity, such that an individual may carry the \textbf{same haplotype on both homologous chromosomes}. Such an individual, being homozygous across the classical MHC loci, expresses fewer distinct allotypes than a fully heterozygous individual---typically one product from each of HLA-A, HLA-B, and HLA-C, together with a correspondingly reduced class II repertoire.

          \myfig{0.35}{immunology_figures/mhc_recomb}{Recombinant HLA gene inheritance; \textit{crossing over produces recombinant MHC haplotype received by offspring}}{fig:mhc_recomb}

      \end{enumerate}

    \end{core}

  \subsection{MHC Antigen Preparation}

    \begin{core}[label=core:endogenous_mhci]{Endogenous MHC I (T$_{\text{c}}$ Cells)}

      \begin{enumerate}[label=(\arabic*)]

        \item \textbf{Proteasomal Degradation of Antigen to Peptide}

          <FIGURE>

          % The antigen is targeted for degradation by binding to ubiquitin—a small protein—and is sent to the interior of a proteasome
          Cytosolic proteins destined for degradation are commonly tagged with chains of the small protein \textbf{ubiquitin}, which direct them to the proteasome for ATP-dependent unfolding and proteolysis.

          % Proteosome is made up of a protein with a 7-fold symmetry (associated with danger) and is more active under conditions of stress and damage by their nature
          The proteolytic 20S core of the proteasome consists of four stacked heptameric rings. Cellular stress and inflammatory signals can increase protein turnover and alter proteasome composition, thereby shaping the pool of peptides available for MHC class I presentation.

          % LMP2 and LMP7 replace proteasomal subunits under inflammatory conditions and promote degradation into the correct peptide size (9 amino acids)
          During inflammation, interferon signaling promotes incorporation of the inducible catalytic subunits LMP2 and LMP7 into the \textbf{immunoproteasome}. These substitutions favor the production of peptides with termini suitable for MHC class I binding, although the products may require additional trimming and are not uniformly nine residues long.

          % There are also complex backup systems for this
          Additional cytosolic peptidases and alternative proteolytic pathways can generate or trim antigenic peptides when conventional proteasomal processing is insufficient.

          % LMP2 and LMP7 genes are located within the MHC II gene sequence and are polymorphic
          LMP2 and LMP7 are encoded by the \textit{PSMB9} and \textit{PSMB8} genes, respectively, which lie within the MHC class II region and exhibit genetic polymorphism.

        \item \textbf{Transport to the RER}

          <FIGURE>

          % The TAP1 and TAP2 bind to the peptides, hydrolyze ATP, and drive them into the RER lumen
          The transporter associated with antigen processing (TAP), a heterodimer of TAP1 and TAP2, binds peptides on the cytosolic face of the rough endoplasmic reticulum (RER) and uses ATP hydrolysis to translocate them into the RER lumen.

          % TAP1 and TAP2 genes are located within the MHC II gene sequence and are polymorphic
          The polymorphic \textit{TAP1} and \textit{TAP2} genes are located within the MHC class II region.

          % TAP Protein Structure
          \textbf{TAP structure} reflects its function as an ATP-binding cassette transporter.

          \begin{itemize}

            % Heterodimer
            \item TAP is an obligate heterodimer composed of TAP1 and TAP2.

            % Cytosolic binding side for peptides
            \item Its peptide-binding site is accessible from the cytosol.

            % ATP hydrolysis site
            \item Cytosolic nucleotide-binding domains hydrolyze ATP to power peptide translocation.

            % Hydrophobic anchoring regions with 3 switchbacks per peptide
            \item Each subunit contains multiple hydrophobic transmembrane helices that together form the peptide-conducting pathway through the RER membrane.

          \end{itemize}

        \item \textbf{Assembly of MHC I}

          \begin{enumerate}

            \item \textit{Calnexin-α to α-β}

              <FIGURE>

              % The α chain is inserted into the membrane, and the β microglobulin chain sent through to the lumen
              The newly synthesized MHC class I α chain is inserted into the RER membrane, while soluble β2-microglobulin enters the RER lumen.

              % Here calnexin ensures the α chain’s proper folding
              The membrane-bound chaperone calnexin associates with the α chain and promotes its proper folding.

              % Once the β microglobulin chain bind to the α chain, calnexin unbinds
              Binding of β2-microglobulin stabilizes the α chain and permits calnexin to dissociate.

            \item \textit{α-β-Calreticulin-Tapasin-TAP-ERp57}

              <FIGURE>

              % The complex associates within calreticulin and tapasin, and then ERp57 binds over the cleft
              The partially assembled MHC class I heterodimer joins calreticulin, tapasin, and ERp57 to form the peptide-loading complex associated with TAP.

              % Calreticulin and ERp57 both help stabilize the molecule and prepare it for peptide binding
              Calreticulin and ERp57 stabilize the peptide-receptive MHC class I molecule and support quality control during peptide loading.

              % Tapasin brings over the TAP complex over to the pre-peptide MHC I
              Tapasin bridges MHC class I to TAP, positioning the empty peptide-binding groove near the source of newly translocated peptides.

            \item \textit{α-β-Calreticulin-Tapasin-TAP-Peptide}

              <FIGURE>

              % TAP protein hands off the peptide via tapasin to the cleft which displaces ERp57
              Peptides released by TAP sample the MHC class I groove, while the peptide-loading complex favors the selection of high-affinity cargo. Stable peptide binding induces conformational maturation and ultimately releases the loaded MHC class I molecule from its chaperones.

            \item \textit{α-β-Peptide to Golgi}

              <FIGURE>

              % Calreticulin and tapasin dissociate from the MHC I-peptide complex and travel to the Golgi
              After calreticulin and tapasin dissociate, the stable peptide--MHC class I complex exits the RER and proceeds through the Golgi apparatus.

          \end{enumerate}

        \item \textbf{Transportation to Exterior}

          <FIGURE>

          % After final sorting the Golgi, the vesicle with MHC I-peptide complex is released by exocytosis onto the membrane exterior
          Following processing and sorting in the Golgi, vesicles deliver peptide--MHC class I complexes to the plasma membrane by exocytosis, exposing the peptide-binding platform to surveillance by CD8$^+$ T cells.

      \end{enumerate}

    \end{core}

    \begin{core}[label=core:cross_presentation]{Exogenous MHC I / Cross Presentation (T$_{\text{c}}$ Cells)}

      \begin{enumerate}[label=(\arabic*)]

        \item \textbf{Antigen Uptake}

          <FIGURE>

          % Through endocytosis, the cell picks up the exogenous antigen
          Specialized antigen-presenting cells, particularly dendritic cells, internalize exogenous antigens through endocytosis or phagocytosis.

        \item \textbf{Proteasomal Degradation of Antigen to Peptide}

          <FIGURE>

          % After getting degraded within the endosome through degradation stages, the peptide fragments travel to and fuse with the RER, exposing the exogenous peptides to MHC I proteins
          Internalized proteins are processed through either a cytosolic route, in which antigen reaches the cytosol for proteasomal degradation and TAP-dependent loading, or a vacuolar route, in which processing and class I loading occur within endocytic compartments. The resulting peptide--MHC class I complexes are transported to the cell surface.

      \end{enumerate}

      % NOTE: This process only happens in APCs, and activated TC cells do not attack the APCs because they don’t express co-stimulatory danger signals that promote an attack on the presenting cell
      Cross-presentation is performed most efficiently by specialized dendritic-cell subsets. During \textbf{cross-priming}, these cells provide peptide--MHC class I together with co-stimulatory and cytokine signals that activate naive CD8$^+$ T cells; presentation in the absence of appropriate co-stimulation may instead promote tolerance.

      % The purpose of cross presentation is to lower the response time of TC cells.
      Cross-presentation enables CD8$^+$ T-cell responses against viruses or tumors that do not directly infect professional antigen-presenting cells, thereby allowing timely cytotoxic immunity to cell-associated exogenous antigens.

    \end{core}

    \begin{core}[label=core:exogenous_mhcii]{Exogenous MHC II}

      \begin{enumerate}[label=(\arabic*)]

        \item \textbf{Antigen Processing}

          \begin{enumerate}

            \item \textit{Antigen Uptake}

              <FIGURE>

              % The antigen is internalized through phagocytosis—a specialized form of endocytosis during which there occurs an engulfment of whole particles with pathogen
              Exogenous antigens enter antigen-presenting cells through receptor-mediated endocytosis, macropinocytosis, or phagocytosis. Phagocytosis is specialized for the engulfment of large particles, including intact microorganisms.

              % The formed vesicle injkj the cell interior is named an endosome
              Endocytosed material enters an endosome, whereas phagocytosed material is initially enclosed within a phagosome.

            \item \textit{Hydrolysis in Vesicles}

              <FIGURE>

              % An endosome goes through three stages (each one more acidic and hydrolytic) and use lysosomes to degrade the foreign antigens (and remove carbohydrates and lipids), cutting them up to their proper size (13-18 amino acids long)
              As endocytic compartments mature, progressive acidification activates resident hydrolases that degrade internalized proteins and remove associated carbohydrates and lipids. Proteolysis generates peptides of variable length, commonly approximately 13--18 residues, that can bind the open-ended groove of MHC class II.

              \begin{itemize}

                % Early endosome, pH 6.0 - 6.5
                \item Early endosomes maintain a mildly acidic pH of approximately 6.0--6.5.

                % Late endosomes of endolysosomes, pH 5.0 - 6.0
                \item Late endosomes and endolysosomes acidify further, generally reaching approximately pH 5.0--6.0.

                % Lysosomes (or phagolysosomes), pH 4.5 - 5.0
                \item Lysosomes and phagolysosomes provide the most acidic and hydrolytic environment, typically near pH 4.5--5.0.

              \end{itemize}

          \end{enumerate}

        \item \textbf{MHC II}

          \begin{enumerate}

            \item \textit{Preparation of MHC II} 

              <FIGURE>

              % After translation of the MHC II αβ heterodimer onto the RER membrane, the heterodimer associates with the invariant chain (Ii)
              Following cotranslational insertion of the MHC class II α and β chains into the RER membrane, the heterodimer associates with the invariant chain (Ii; CD74).

              % The figure only shows one invariant chain binding, but in reality, the Ii comes as a trimer—three MHC II heterodimers bind at a time
              Although a simplified diagram may depict a single invariant chain, Ii trimerizes and can associate with three MHC class II αβ heterodimers.

              % The invariant chain [1] blocks the cleft (to prevent immature binding), [2] ensures proper conformation, [3] helps RER exiting, and [4] signals for fusing with phagolysosomes
              The invariant chain occupies the peptide-binding groove to prevent premature loading of RER-derived peptides, stabilizes the nascent heterodimer, promotes exit from the RER, and contains sorting signals that direct the complex into the endosomal pathway.

            \item \textit{Transportation of Golgi}

              <FIGURE>

              % The αβ-Ii complex is then transported to the Golgi where it awaits fusion with a phagolysosome
              The MHC class II--Ii complex traverses the Golgi and is subsequently sorted into acidified endosomal compartments containing processed exogenous antigen.

          \end{enumerate}

        \item \textbf{Fusion with Endocytic Vesicles/Phagolysosomes}

          <FIGURE>

          % The αβ-Ii complex fuses with early endosomes after leaving the Golgi
          After leaving the Golgi, vesicles containing MHC class II--Ii complexes enter the endocytic pathway and mature toward late endosomal peptide-loading compartments.

        \item \textbf{Invariant Chain Degradation}

          <FIGURE>

          % After the development of the endosomes into more acidic internal environments, Ii gets degraded, leaving only the CLIP attached to the MHC II and still blocking the cleft region
          Acid-activated proteases progressively degrade the invariant chain, leaving the class II-associated invariant-chain peptide (CLIP) within the MHC class II groove.

        \item \textbf{CLIP Removal}

          <FIGURE>

          % At the end of antigen degradation into peptides, non-classical MHC (HLA-DM) removes the CLIP, allowing for peptide binding and the formation of a peptide-presenting MHC II protein (which is very stable)
          The nonclassical MHC class II molecule HLA-DM catalyzes CLIP release and edits the peptide repertoire by favoring ligands that form stable complexes with MHC class II.

        \item \textbf{Transportation to Exterior}

          <FIGURE>

          % The peptide-MHC II complex is then transported to the cell surface to be displayed
          Stable peptide--MHC class II complexes are transported to the plasma membrane for recognition by CD4$^+$ T cells.

      \end{enumerate}

    \end{core}

    \begin{core}[label=core:exogenous_cd1mhc]{Exogenous CD1 MHC}

      \begin{enumerate}[label=(\arabic*)]

        % \item \textbf{Glycolipid Uptake}
        \item \textbf{Glycolipid Uptake}: Antigen-presenting cells internalize microbial or endogenous glycolipids into endosomal compartments.

        % \item \textbf{Preparation and Expression of CD1 MHC}
        \item \textbf{Preparation and Expression of CD1}: CD1 molecules assemble in the RER with β2-microglobulin and traffic through the secretory pathway to the plasma membrane.

        % \item \textbf{Retreat of CD1 MHC}
        \item \textbf{Internalization of CD1}: Surface CD1 molecules can be internalized and routed into endosomal compartments through signals in their cytoplasmic tails.

        % \item \textbf{Binding of CD1 MHC to Glycolipid}
        \item \textbf{Glycolipid Loading}: Within the endosomal system, lipid-transfer proteins facilitate the loading of glycolipid antigens into the hydrophobic binding groove of CD1.

        % \item \textbf{Transportation to Exterior}
        \item \textbf{Transport to the Cell Surface}: Loaded glycolipid--CD1 complexes return to the plasma membrane for recognition by lipid-reactive T cells.

      \end{enumerate}

      % NOTE: Several different types of cells and receptors can bind to and respond to the CD1 lipid presentation
      Several lymphocyte populations can recognize or respond to lipid antigens presented by members of the CD1 family.

      \begin{itemize}

        % T cells bind to the lipid-CD1 MHC complex through the γδ receptor (with no CD4 or CD8)
        \item Some γδ T cells recognize lipid--CD1 complexes without requiring CD4 or CD8 co-receptors.

        % TH cells can also bind through the αβ receptors using CD4
        \item Certain conventional αβ T cells, including CD4$^+$ populations, recognize antigens presented by CD1 molecules.

        % NK cells
        \item Natural killer cells may respond indirectly to cytokines and cellular changes associated with CD1-restricted immune responses, although they do not use clonotypic TCRs to recognize lipid--CD1 complexes.

        % invariant natural killer T cells (iNKT) with CD1
        \item Invariant natural killer T cells use a semi-invariant αβ TCR to recognize glycolipids presented specifically by CD1d.

      \end{itemize}

    \end{core}
 
\section{T Cells}

  \subsection{T Cell Receptor}

    \subsubsection{TCR Complex Characteristics}
    \label{sec:tcr}

      <FIGURE>

      % It is important to keep in mind that T cells respond to antigen only when it is presented by MHC molecules encoded by alleles that they themselves possess.

      \begin{core}{αβ and γδ Receptors}

        <FIGURE>

        % T cells can make either the αβ or γδ TCR—a heterodimer (most T cells make the αβ version)
        Each mature T cell expresses either an αβ or a γδ T-cell receptor (TCR), both of which are heterodimers. Most circulating human T cells express the αβ form.

        % Both resemble a truncated Ig and is part of the Ig superfamily with three hypervariable loops in each chain (CDRs)
        Both receptor classes belong to the immunoglobulin superfamily and resemble membrane-bound antigen-binding fragments. Each chain contributes three complementarity-determining regions (CDRs) to the antigen-recognition surface.

        % The positively charged transmembrane regions of the TCR and negatively charged transmembrane regions of the CD3 signal transducers form a salt bridge attraction that holds the complex together
        Basic residues within the TCR transmembrane segments form ionic interactions with acidic residues in the CD3 signaling subunits, stabilizing assembly of the complete receptor complex.

        % During development, the TCR does not go through affinity maturation like B cells, thus they need the supplement receptors (CD4/8) in order to strengthen this less specific binding
        Unlike immunoglobulins, TCRs do not undergo somatic hypermutation or affinity maturation after activation. The CD4 or CD8 co-receptor stabilizes recognition of the appropriate MHC class and recruits signaling machinery to the engaged TCR.

        % αβ Receptor
        \textbf{αβ receptors} undergo thymic selection to generate a repertoire capable of recognizing foreign peptide in the context of self MHC while remaining tolerant of self antigens.

        % Undergoes selection to produce a highly specific recognition molecule
        The resulting receptor provides highly specific recognition of a peptide--MHC complex.

        % γδ Receptor
        \textbf{γδ receptors} support a more innate-like mode of immune surveillance, although their antigen recognition remains clonally distributed and can be highly specific.

        % Undergoes selection to produce a quasi innate recognition molecule (some specificity but not as much as αβ)
        Many γδ T cells respond rapidly to conserved stress-associated or microbial ligands without conventional peptide--MHC restriction.

        % Mostly for recognizing lipid antigens of bacteria and parasites
        Depending on the subset, γδ T cells can recognize microbial metabolites, lipids, and stress-induced host molecules associated with bacterial or parasitic infection.

        % Important in patrolling mucosa and epithelium
        These cells are particularly enriched at mucosal and epithelial barriers, where they provide rapid tissue surveillance.

      \end{core}

      \begin{core}{γε and δε Heterodimers}

        <FIGURE>

        % All chains (γ, δ, ε) are part of the Ig superfamily
        The CD3γ, CD3δ, and CD3ε chains possess extracellular immunoglobulin-superfamily domains and assemble as CD3γε and CD3δε heterodimers.

        % All peptides have an immunoreceptor tyrosine-based activation motif (ITAM) at the cell-interior side of the receptor
        Each CD3 chain contains one immunoreceptor tyrosine-based activation motif (ITAM) within its cytoplasmic tail, allowing ligand engagement to be converted into an intracellular signal.

      \end{core}

      \begin{core}{ζζ or ζη Heterodimers}

        <FIGURE>

        % Chains are NOT part of the Ig superfamily
        The ζ and η chains do not contain extracellular immunoglobulin-superfamily domains.

        % Only 9 amino acids to the exterior and a long extended cytoplasmic tail
        They have extremely short extracellular segments and long cytoplasmic tails specialized for signaling.

        % Has 3 ITAMs rather can one
        Each ζ-family chain contains three ITAMs rather than the single ITAM present in each CD3γ, CD3δ, or CD3ε chain. The predominant TCR complex contains a ζζ homodimer, whereas ζη heterodimers are less common.

      \end{core}

    \subsubsection{CD4 \& CD8}

      \begin{core}{CD4}

        <FIGURE>

        % Used by TH cells to bind to MHC II
        CD4 is a co-receptor expressed by helper T cells and other subsets that binds nonpolymorphic regions of MHC class II during TCR recognition.

        % Monomer
        CD4 functions as a monomeric transmembrane glycoprotein.

        % 4 extracellular Ig domains
        Its extracellular region contains four immunoglobulin-superfamily domains.

        % N-terminal binds to the base of MHC II (junction of α2 and β2 domains) at a hydrophobic pocket
        The amino-terminal D1 domain contacts a nonpolymorphic surface formed principally by the membrane-proximal β2 domain of MHC class II, away from the peptide-binding groove.

        % C-terminal possess 3 serine regions that can be phosphorylated
        The cytoplasmic tail associates noncovalently with the Src-family kinase Lck, thereby positioning Lck near the TCR--CD3 complex during antigen recognition.

      \end{core}

      \begin{core}{CD8}

        <FIGURE>

        % Used by TC cells to bind to MHC I
        CD8 is the characteristic co-receptor of cytotoxic T cells and binds a nonpolymorphic region of MHC class I.

        % Dimer—α-β hetero or α-α homo linked by disulfide
        It is expressed as a disulfide-linked CD8αβ heterodimer or CD8αα homodimer.

        % 1 extracellular Ig domain per peptide
        Each CD8 chain contributes one extracellular immunoglobulin-superfamily domain.

        % N-terminal binds to the α2 and β microglobulin
        The amino-terminal domains bind primarily to the α3 domain of the MHC class I α chain, with additional contacts involving β2-microglobulin.

        % C-terminal possess amino acid regions that can be phosphorylated
        The cytoplasmic tail of CD8α associates with Lck, coupling co-receptor engagement to phosphorylation of the TCR signaling complex.

      \end{core}

    \subsubsection{TCR Genes and Rearragement}

      \begin{core}{Gene Sequence and Receptor Overview}

        <FIGURE>

        \begin{enumerate}

          \item \textbf{α \& δ Peptide} 

            <FIGURE>

            % Chromosome 14 Long Arm
            The TCR α and δ loci occupy a shared region on the long arm of chromosome 14, with the δ locus nested within the α locus.

            % Alpha (α): Resembles the light chain Ig gene because of absence of D regions
            % ~40 V --- N/A D --- ~50 J --- 1 C
            Like an immunoglobulin light-chain locus, the α locus lacks D segments and assembles a VJ exon from approximately 40 functional V segments, approximately 50 J segments, and one C region.

            % Delta (δ): Resembles the heavy chain Ig gene because of presence of D regions
            % ~3 (or more) V --- ~3 D --- ~4 J --- 2 C
            Like an immunoglobulin heavy-chain locus, the δ locus contains D segments and forms a V(D)J exon. It uses a limited set of V, D, and J segments followed by a single Cδ gene; some V segments are shared with the α locus.

          \item \textbf{β Peptide}

            <FIGURE>

            % Chromosome 7 Long Arm
            The TCR β locus lies on the long arm of chromosome 7.

            % Resembles the heavy chain Ig gene because of presence of D regions
            % ~40 V --- 2 (functional) D --- 6 and 7 J --- 2 C
            Its organization resembles that of an immunoglobulin heavy-chain locus because it includes D segments. Approximately 40 functional V segments precede two D--J--C clusters, whose J clusters contain six and seven functional segments, respectively.

          \item \textbf{γ Peptide}

            <FIGURE>

            % Chromosome 7 Short Arm
            The TCR γ locus lies on the short arm of chromosome 7.

            % Resembles the light chain Ig gene because of absence of D regions
            % ~5 (functional) V --- N/A D --- ~5 J --- 2 C
            Its organization resembles that of an immunoglobulin light-chain locus because it lacks D segments. A limited set of functional V segments rearranges with J segments associated with two principal Cγ regions.

        \end{enumerate}

      \end{core}

      \begin{core}{α Rearrangement}

        \begin{enumerate}[label=(\arabic*)]

          \item \textbf{Gene Sequence}

            <FIGURE>

          \item \textbf{δ Removal \& VL-J Joining}

            <FIGURE>

            % **Splicing and joining of the sequences are carried out through the same process of RSS manipulation by RAG enzymes done in Ig gene rearrangement ALSO P and N nucleotide addition occurs and further adds to the variability
            RAG1 and RAG2 recognize recombination signal sequences and mediate V--J joining through the same V(D)J recombination mechanism used at immunoglobulin loci. Imprecise junction formation, including P- and N-nucleotide addition, further increases receptor diversity. Rearrangement of the α locus deletes the intervening δ locus from that chromosome.

          \item \textbf{Result}

            <FIGURE>

            % [Variable \# of VL] - [VL] - [J] - [Variable \# of J] - [Intron] - [C]
            A completed α-chain locus contains a selected V segment joined to one J segment, followed downstream by unrearranged J segments, an intron, and the Cα region.

            % Next is EXPRESSION
            Transcription and RNA processing subsequently produce the mature α-chain message.

        \end{enumerate}

      \end{core}

      \begin{core}{β Rearrangement}

        \begin{enumerate}[label=(\arabic*)]

          \item \textbf{Gene Sequence}

            <FIGURE>

          % \item \textbf{VL-D/J Joining}
          \item \textbf{V--D--J Joining}

            \begin{itemize}

              % \item \textit{VL-DJ}
              \item \textit{V--D--J rearrangement}

                <FIGURE>

                % VL joins with a D
                The TCR β locus rearranges sequentially: a D segment first joins a J segment, after which a V segment joins the resulting DJ unit.

                % **Splicing and joining of the sequences are carried out through the same process of RSS manipulation by RAG enzymes done in Ig gene rearrangement ALSO P and N nucleotide addition occurs and further adds to the variability
                RAG-mediated cleavage at recombination signal sequences and imprecise end joining generate the rearranged VDJ exon. P- and N-nucleotide addition at both junctions substantially increases CDR3 diversity.

              % \item \textit{VL-J}
              \item \textit{Direct V--J joining}

                <FIGURE>

                % VL joins with a J (all Ds get cut out)
                Direct V--J joining is not a normal productive pathway at the TCR β locus because the 12/23 recombination constraints enforce incorporation of a Dβ segment.

                % **Splicing and joining of the sequences are carried out through the same process of RSS manipulation by RAG enzymes done in Ig gene rearrangement ALSO P and N nucleotide addition occurs and further adds to the variability
                Consequently, productive β-chain rearrangement uses V--D--J recombination rather than the alternative V--J scheme depicted in the rough notes.

            \end{itemize}

          \item \textbf{Result}

            \begin{itemize}

              % \item \textit{VL-DJ}
              \item \textit{V--D--J product}

                <FIGURE>

              % \item \textit{VL-J}
              \item \textit{Noncanonical V--J product}

                <FIGURE>

            \end{itemize}

        \end{enumerate}

        % Next occurs EXPRESSION
        A productive VDJ exon is subsequently transcribed and spliced to the appropriate Cβ region for β-chain expression.

      \end{core}

      \begin{core}{γ Rearrangement}

        \begin{enumerate}[label=(\arabic*)]

          \item \textbf{Gene Sequence}

            <FIGURE>

          % \item \textbf{VL-D/J Joining}
          \item \textbf{V--J Joining}

            <FIGURE>

            % **Splicing and joining of the sequences are carried out through the same process of RSS manipulation by RAG enzymes done in Ig gene rearrangement ALSO P and N nucleotide addition occurs and further adds to the variability
            RAG1 and RAG2 mediate V--J recombination at the γ locus through recombination signal sequences. Junctional processing, including P- and N-nucleotide addition, diversifies the resulting CDR3 sequence.

          \item \textbf{Result}

            <FIGURE>

        \end{enumerate}

        % Next occurs EXPRESSION
        The rearranged VJ exon is subsequently transcribed and spliced to a compatible Cγ region for receptor expression.

      \end{core}

      \begin{core}{δ Rearrangement}

        \begin{enumerate}[label=(\arabic*)]

          \item \textbf{Gene Sequence}

            <FIGURE>

          % \item \textbf{VL-D/J Joining \& Result}
          \item \textbf{V--D--J Joining and Result}

            <FIGURE>

            % The VL sequence either joins to a D or J depending on what it splices out
            Productive TCR δ rearrangement generally joins a V segment to one or more D segments and then to a J segment, creating a V(D)J exon.

            % NOTE: This rearrangement can add more than one D to the sequence as well
            Unlike most antigen-receptor loci, the arrangement of recombination signal sequences at the δ locus can permit incorporation of multiple D segments into a single rearranged chain.

            % **Splicing and joining of the sequences are carried out through the same process of RSS manipulation by RAG enzymes done in Ig gene rearrangement ALSO P and N nucleotide addition occurs and further adds to the variability
            RAG-mediated recombination and junctional processing, including P- and N-nucleotide addition, generate extensive diversity within the δ-chain CDR3 region.

        \end{enumerate}

        % Next occurs EXPRESSION
        Transcription and RNA processing then join the rearranged variable exon to Cδ for expression of the mature δ chain.

      \end{core}

      \begin{core}{Subsequent Expression}

        \begin{enumerate}[label=(\arabic*)]

          % \item The message gets trancribed
          \item The rearranged locus is transcribed into a primary RNA transcript.

            % NOTE: TCR C regions always has its bulk C gene sequence followed by an extracellular gene (XC), a transmembrane gene (TM), and a cytoplasmic tail gene (CA)
            TCR constant regions are encoded by exons that specify the extracellular constant domain, connecting peptide, transmembrane segment, and short cytoplasmic tail.

            <FIGURE>

          % \item The message precursor is capped and gets a poly-A-tail
          \item The primary transcript receives a 5$'$ cap and a 3$'$ poly(A) tail.

          % \item Processing occurs (removing the introns and other extra Js)
          \item RNA splicing removes introns and excludes unused downstream gene segments, joining the rearranged variable exon to the appropriate constant-region exons.

            <FIGURE>

          % \item Translation of the L region attaches the ribosome to the RER and inserts into the lumen with the C-terminal bound to the ER membrane
          \item Translation of the amino-terminal leader sequence directs the ribosome to the RER, where the nascent receptor chain enters the secretory pathway and its transmembrane segment anchors it in the ER membrane.

          % \item The L region is clipped off
          \item Signal peptidase removes the leader peptide during translocation into the RER.

          % \item Enzymes in the Golgi form disulfide bonds
          \item Folding, disulfide-bond formation, and initial glycosylation occur primarily in the RER, followed by additional glycan processing in the Golgi.

          % \item Vesicles with finished TCRs get transported to and fuses with the plasma membrane
          \item Vesicles transport fully assembled TCR--CD3 complexes from the Golgi to the plasma membrane, where vesicle fusion displays the receptors on the cell surface.

        \end{enumerate}

      \end{core}

      \begin{core}{History: \textit{Hendrick and Davis}}

        % Hendrick and Davis were responsible for the gene identification of TCRs
        Mark Davis and Stephen Hedrick helped identify genes encoding the TCR through experiments designed to isolate T-cell-specific, somatically rearranged transcripts.

        \begin{enumerate}[label=(\arabic*)]

          \item \textbf{Selecting Membrane Bound Polysomes}

            <FIGURE>

            % H\&D first separated membrane bound from cytoplasmic proteins by isolating mRNA from membrane bound polysomes
            Hedrick and Davis enriched for transcripts encoding secreted or membrane proteins by isolating mRNA associated with membrane-bound polysomes.

            % T cells—TCR, CD3, CD4/8
            In T cells, this fraction included transcripts encoding the TCR, CD3 subunits, and the CD4 or CD8 co-receptors.

            % B cells—Ig receptor (BjkjkCR), co-receptors
            In B cells, it included transcripts encoding membrane immunoglobulin and associated co-receptors.

            % Both of these cells with have other things in common as well
            Because T and B lymphocytes also express many shared membrane proteins, further subtraction was required to isolate T-cell-specific candidates.

          \item \textbf{Subtractive Hybridization}

            <FIGURE>

            % H\&D used membrane bound B cell genes to selectively hybridize with and remove any overlapping radioactive cDNA from T Cells, leaving only the radioactive cDNA unique to T cells behind
            During subtractive hybridization, B-cell-derived nucleic acids were used to remove cDNAs shared by the two lymphocyte lineages, enriching the remaining library for transcripts preferentially expressed by T cells.

          \item \textbf{Gene Comparison Across Multiple T Cells}

            <FIGURE>

            % Now there are only 10 different cDNA clones remained (unique to T cells) and we have to find which ones are rearranged (which would be the TCR)
            The enriched candidate pool contained approximately ten T-cell-specific cDNA clones, which were then screened for the somatic DNA rearrangement expected of an antigen-receptor gene.

            % We run each cDNA portion across multiple T cells lineages
            Each candidate probe was compared across independently derived T-cell clones and a non-T-cell control by genomic blotting.

            % If the cDNA portion (for example, g1) is not rearranged, the length traveled in the gel electrophoresis would be the same across all T cell lineages and B cells (control) as well
            A candidate that remained in germline configuration produced the same restriction-fragment pattern across the different cell lineages.

            % If the cDNA portion (for example, g2) is rearranged, the length traveled in the gel electrophoresis would be different across different T cell lineages because of their variable lengths
            By contrast, a rearranged candidate produced clone-specific restriction fragments because independent T cells had assembled different variable-region sequences.

            % The rearranged gene were 2 out of the 10 cDNAs, thus locating the genes for the TCR
            Two candidates exhibited T-cell-specific rearrangement, enabling identification of genes encoding the TCR chains.

        \end{enumerate}

        % Alloreactivity is T cells being able to bind with foreign MHCs, which it should not be able to do (Figure on MHC Characteristics)
        \textbf{Alloreactivity} is the ability of a T cell selected on self MHC to recognize a structurally distinct allogeneic MHC--peptide complex. This cross-reactivity underlies strong T-cell responses to genetically mismatched transplants (see the figure on MHC characteristics).

      \end{core}

  \subsection{T Cell Development}

    \begin{core}{Developmental Pathway}

      <TABLE>

    \end{core}

  \subsection{TCR Antigen Response}

    \begin{core}{Antigen Response Pathway}

      \begin{enumerate}[label=(\arabic*)]

        \item \textbf{Antigen Recognition}

          % APC presents an antigenic peptide on MHC II
          A professional antigen-presenting cell displays an antigen-derived peptide within the binding groove of an MHC class II molecule.

        \item \textbf{Binding}

          % TH cell binds to this complex through TCR–CD3 complex, aided by CD4
          A cognate CD4$^+$ helper T cell binds the peptide--MHC class II complex through its TCR, while CD4 stabilizes the interaction and recruits Lck to the CD3 signaling apparatus.

          % *Superantigens are viral or bacterial proteins that nail together TCR-antigen-MHC complexes, hyperstimulating the TH cells and overproducing cytokines (which could kill you)
          \textit{Superantigens} are microbial proteins that crosslink particular TCR Vβ regions to MHC class II outside the conventional peptide-binding interface. This noncanonical interaction activates unusually large T-cell populations and can provoke life-threatening systemic cytokine release.

        \item \textbf{Signaling Cascade}

          % A signaling cascade starts towards the nucleus and pushes the cell into cell division
          TCR engagement and co-stimulation initiate interconnected signaling cascades that reach the nucleus, alter gene expression, and commit the cell to activation, proliferation, and differentiation.

        \item \textbf{Transcriptional Activation}

          \begin{itemize}

            \item \textit{Immediate Genes}

              % Genes for transcription factors for inflammatory response
              Immediate signaling activates pre-existing transcription factors that coordinate the early inflammatory and proliferative response.

              % —NFAT, AP-1, NF-κB
              Principal factors include NFAT, AP-1, and NF-κB.

            \item \textit{Early Genes}

              % Cytokines, Interleukins and their receptors
              Early-response genes encode cytokines, cytokine receptors, and regulatory proteins that reinforce activation and determine developmental fate.

              % *This is where developmental branching occurs, and the TH cells differentiate into either TH1, TH2, THreg, or TH17 (sometimes referred to as TH3)
              At this stage, naive helper T cells can differentiate into specialized subsets, including T$_{\text{H}}$1, T$_{\text{H}}$2, T$_{\text{H}}$17, and regulatory T cells. T$_{\text{H}}$17 cells should not be designated T$_{\text{H}}$3, a term historically used for a distinct regulatory phenotype.

              % Inducing cytokines are produced by other immune cells like macrophages and DCs which during the production of master transcription factors
              Cytokines produced by dendritic cells, macrophages, and other cells activate lineage-defining transcription factors in the responding T cell.

              % —[Type]: [inducing cytokine] → [master transcription factor]
              The principal relationships between polarizing cytokines and lineage-defining transcription factors include the following:

              \begin{itemize}

                % —Th1: IL-12 → T-bet
                \item IL-12 promotes T$_{\text{H}}$1 differentiation through T-bet;

                % —Th2: IL-4 → GATA3
                \item IL-4 promotes T$_{\text{H}}$2 differentiation through GATA3;

                % —Th17: TGF-β + IL-6 → RORγt
                \item TGF-β together with inflammatory cytokines such as IL-6 promotes T$_{\text{H}}$17 differentiation through RORγt; and

                % —Treg: TGF-β + IL-2 → FOXP3
                \item TGF-β and IL-2 promote peripheral regulatory T-cell differentiation and FOXP3 expression.

              \end{itemize}

            \item \textit{Late Genes}

              % Cell adhesion molecules (CAMs) and effector cytokines
              Late-response genes encode effector cytokines, adhesion molecules, and other proteins required for migration and specialized immune function.

          \end{itemize}

      \end{enumerate}

    \end{core}

    \begin{core}{TCR Antigen Recognition to ZAP70}

      <FIGURE>

      \begin{enumerate}[label=(\arabic*)]

        % \item APCs present their MHC II and antigen complex to the TCR
        \item An antigen-presenting cell displays a peptide--MHC class II complex to a cognate TCR.

        % \item CD4 helps strengthen this binding
        \item CD4 binds the MHC class II molecule, stabilizes the receptor complex, and positions its associated kinase Lck near CD3.

        % \item p56lck (which is associated with CD4/8 on its cytoplasmic tail) phosphorylates the ITAMs, providing a docking site for ZAP 70
        \item Co-receptor-associated p56$^{\text{Lck}}$ phosphorylates tyrosines within CD3 and ζ-chain ITAMs, creating docking sites for ZAP-70.

        % \item ZAP 70 docks on the ITAMs and get phosphorylated by p56lck and activated
        \item ZAP-70 binds doubly phosphorylated ITAMs through its tandem SH2 domains and is phosphorylated and activated by Lck.

        % \item ZAP-70 goes on to phosphorylate a variety of molecules, setting off other downstream pathways through transcriptional activation
        \item Activated ZAP-70 phosphorylates adaptor proteins, including LAT and SLP-76, which organize downstream pathways leading to transcriptional activation.

        \begin{itemize}

          % —DAG-IP3 formed through hydrolysis of membrane phospholipid by phospholipase C, causes ER to release Ca2+
          \item PLC-γ1 hydrolyzes PIP$_2$ into diacylglycerol (DAG) and inositol 1,4,5-trisphosphate (IP$_3$); IP$_3$ then stimulates Ca$^{2+}$ release from the RER.

          % —Ca2+ activates NFAT (up regulates immune response)
          \item Ca$^{2+}$ binds calmodulin and activates calcineurin, which dephosphorylates NFAT and permits its entry into the nucleus.

          % —DAG activates NF-κB (up regulates immune response)
          \item DAG activates protein kinase C-θ and contributes to activation of the canonical NF-κB pathway.

          % —Activation of Ras, the G protein
          \item DAG-dependent signaling also activates the small GTPase Ras.

          % —Transcription factors activated by MAP kinases turn on genes that trigger cell division
          \item Ras activates mitogen-activated protein kinase cascades and AP-1, inducing genes required for proliferation and effector differentiation.

        \end{itemize}

      \end{enumerate}

      % Regulation
      \textbf{Co-stimulatory regulation} determines whether TCR signaling produces productive activation or functional unresponsiveness.

      % B7 is a membrane bound protein on APCs binding either to the CD28 or CTLA-4 of a T cell
      The B7 ligands CD80 and CD86 are membrane proteins on antigen-presenting cells that bind either the activating receptor CD28 or the inhibitory receptor CTLA-4 on T cells.

      % CD28 Binding (Upregulation)
      \textbf{CD28 engagement} amplifies TCR signaling and promotes cell survival, metabolism, cytokine production, and proliferation.

      % The cell will begin proliferating and dividing, transcribing the message for both IL-2 and the high affinity version of its TCR α chain and enter into a positive feedback loop.
      CD28 co-stimulation increases transcription of IL-2 and expression of the high-affinity IL-2 receptor α chain (CD25), establishing an autocrine positive-feedback loop that supports clonal expansion. It does not generate a higher-affinity TCR α chain.

      % *An APC without the B7 for CD28 would turn the T cell off
      Recognition of antigen by a naive T cell in the absence of adequate CD28 co-stimulation can induce anergy, deletion, or another form of peripheral tolerance.

      % CTLA-4 Binding (Downregulation)
      \textbf{CTLA-4 engagement} opposes CD28-mediated co-stimulation and limits the magnitude of T-cell activation.

      % The T cell will being to downregulate its development
      CTLA-4 therefore restrains proliferation and effector development after activation.

      % *Resting T cells initially have no CTLA-4, but begin to produce it during activation
      Conventional resting T cells express little surface CTLA-4, but activation induces its expression and trafficking to the plasma membrane; regulatory T cells express CTLA-4 constitutively.

    \end{core}

    \begin{core}{Lipid Rafts}

      % Over the course of the signaling lipid rafts migrate towards the outside of the synapse region as a result of CAMs associating CD45 to CD22
      During immunological-synapse formation, TCR microclusters and associated lipid-ordered membrane domains reorganize dynamically. Large phosphatases such as CD45 are initially excluded from close-contact signaling zones, while adhesion molecules and the cytoskeleton segregate receptors into central and peripheral supramolecular activation clusters; CD22 is a B-cell protein and does not organize the T-cell synapse.

      \begin{enumerate}[label=(\arabic*)]

        % \item TCR, lipid raft after antigen binding
        \item \textbf{Initial engagement}: TCR microclusters form rapidly at sites of peptide--MHC binding and recruit proximal signaling proteins.

          <FIGURE>

        % \item Synapse during mid-activation
        \item \textbf{Intermediate organization}: Actin-dependent transport moves signaling microclusters inward while integrins establish a surrounding adhesive zone.

          <FIGURE>

        % \item Synapse, mature (bull’s eye formation)
        \item \textbf{Mature synapse}: The interface develops a bull's-eye organization, with a central supramolecular activation cluster surrounded by a peripheral adhesion-rich ring.

          <FIGURE>

      \end{enumerate}

    \end{core}
  
\section{Cytokines and Signals}

  \subsection{Receptor Signaling Types}

    \begin{core}{Immunoglobulin}

      <FIGURE>

      % Structure
      \textbf{Structure}: Immunoglobulin-superfamily receptors share extracellular domains built from antiparallel β sheets arranged in a characteristic immunoglobulin fold.

      % Receptors have β pleated sheet (Ig) domains and membrane span characteristics of family
      Individual family members combine these domains with distinct transmembrane and cytoplasmic signaling arrangements suited to their functions.

      % Ligands—antigens, integrins, growth factors
      \textbf{Ligands}: Members of this broad family participate in antigen recognition, cell adhesion, and growth-factor signaling and therefore interact with diverse ligands rather than a single ligand class.

    \end{core}

    \begin{core}{Type I Receptors}

      <FIGURE>

      % Structure
      \textbf{Structure}: Type I cytokine receptors are single-pass membrane proteins that commonly assemble as heterodimers or higher-order complexes.

      % Heterodimers or trimers with interchangeable β and γ subunits across different receptors
      Several receptors share signaling subunits, such as the common β chain or common γ chain, allowing related cytokines to activate overlapping pathways.

      % Composed of fibronectin domain (like Ig domains), two β pleated sheets running opposite
      Their extracellular cytokine-receptor homology domains are formed primarily from antiparallel β strands arranged in fibronectin type III-like folds.

      % β,γ N terminals—Four conserved cysteines that stabilize the structure using 2 disulfide bonds
      The membrane-distal portion contains four conserved cysteine residues that form stabilizing disulfide bonds.

      % β,γ proximal domains—conserved [tryptophan]-[serine]-[whatever]-[tryptophan]-[serine] sequence (function unknown)
      A conserved membrane-proximal WSXWS motif contributes to receptor folding, trafficking, and ligand-responsive organization.

      % β,γ C terminals—domains have sites that recruit cytosolic tyrosine kinases (phosphorylation)
      Because these receptors lack intrinsic kinase activity, their cytoplasmic tails recruit Janus kinases (JAKs) that phosphorylate receptor-associated signaling proteins.

      % Ligands—a small monomer with 4 α helical domains
      \textbf{Ligands}: Many type I cytokines are small monomeric proteins organized as four-α-helix bundles.

      % JAK-STAT Pathway—
      \textbf{JAK--STAT pathway}: Ligand-dependent receptor assembly initiates a direct route from the plasma membrane to transcriptional regulation.

      \begin{enumerate}[label=(\arabic*)]

        % Ligand binding brings JAKs together and activates
        \item Ligand binding changes receptor geometry and juxtaposes receptor-associated JAKs, permitting their activation.

        % P-JAKs phosphorylate the cytoplasmic receptor tails (docking site for STAT)
        \item Activated JAKs phosphorylate tyrosine residues on the receptor tails, creating docking sites for STAT proteins.

        % JAKs phosphorylate the STATs, which migrate to the nucleus, and up-regulate other genes
        \item Recruited STATs are phosphorylated, dimerize, enter the nucleus, and regulate expression of cytokine-responsive genes.

      \end{enumerate}

    \end{core}

    \begin{core}{Type II Receptors}

      <FIGURE>

      % Structure
      \textbf{Structure}: Type II cytokine receptors are single-pass membrane proteins that assemble as homo- or heterodimeric signaling complexes.

      % β, γ heterodimers with extensive cytoplasmic signaling domains
      Their subunits possess cytoplasmic regions that associate with JAK-family kinases and recruit downstream signaling proteins.

      % Fibronectin domain
      The extracellular portions contain fibronectin type III-like cytokine-receptor domains.

      % β,γ N terminals—CCCC that stabilize the structure using 2 disulfide bonds
      Conserved cysteine residues form disulfide bonds that stabilize the extracellular structure, but type II receptors lack the canonical WSXWS motif of type I receptors.

      % Ligands—series of α helices
      \textbf{Ligands}: Type II receptors bind helical cytokines, principally interferons and members of the IL-10 family.

      % JAK-STAT Pathway—SAME as Type I
      \textbf{JAK--STAT pathway}: Signal transduction follows the same general logic as in type I cytokine receptors.

      % Ligand binding brings JAKs together and activates
      Ligand binding reorganizes receptor-associated JAKs and promotes their activation.

      % P-JAKs phosphorylate the cytoplasmic receptor tails (docking site for STAT)
      The activated JAKs phosphorylate receptor tails to generate STAT docking sites.

      % JAKs phosphorylate the STATs, which migrate to the nucleus, and up-regulate other genes
      Recruited STATs are phosphorylated, dimerize, and enter the nucleus to regulate target-gene transcription.

      % ***JAKs and STATs have variety (4 and 6 known respectively) and the signaling process uses more than one version of each
      Mammals express four JAK proteins and seven STAT proteins. Individual receptor complexes recruit characteristic combinations of these signaling molecules, allowing different cytokines to produce distinct transcriptional programs.

    \end{core}

    \begin{core}{Interleukin 17}

      <FIGURE>

      % Structure
      \textbf{Structure}: IL-17 receptors constitute a distinct cytokine-receptor family whose subunits can form homo- or heteromeric complexes.

      % Heterodimers, Homodimers, Trimers
      Ligand binding promotes higher-order receptor assembly, although the precise stoichiometry can vary among IL-17 family members.

      % IL17R composed on 2 IL17RAs and 1 IL17RC
      The functional receptor for IL-17A and IL-17F contains IL-17RA and IL-17RC subunits; simplified trimeric descriptions should not be treated as a universal fixed stoichiometry.

      % Ligands—IL-17A: β pleated sheets held together by 2 S=S bonds constructed by CCCC (pro-inflammatory)
      \textbf{Ligands}: IL-17A is a disulfide-linked homodimeric cytokine with a predominantly β-sheet structure and potent pro-inflammatory activity.

      % Activates the Act1 Pathway (NF-κB, MAPK, CCAAT)
      Receptor engagement recruits the adaptor Act1 and activates NF-κB, MAPK, and C/EBP-family transcriptional pathways.

      % **IL-17 is secreted by Th17, iNKT, and γδ T cells (mainly for border patrol)
      IL-17 is produced by T$_{\text{H}}$17 cells as well as several innate-like lymphocyte populations, including γδ T cells and some invariant natural killer T cells, and is particularly important for antimicrobial defense at epithelial barriers.

    \end{core}

    \begin{core}[label=core:tnf]{Tumor Necrosis Factor}

      <FIGURE>

      % Structure
      \textbf{Structure}: Tumor necrosis factor receptors are type I transmembrane proteins characterized by extracellular cysteine-rich domains stabilized by disulfide bonds.

      % Four repeated β pleated sheet domains stabilized by CCCCs’ S=S bonds
      These repeated cysteine-rich modules form the elongated extracellular ligand-binding region.

      % Upon binding TNF, subunits cluster in trimer, leading to activation
      Binding of trimeric TNF ligands promotes productive clustering and conformational rearrangement of receptor subunits.

      % Ligands—juxtacrine factors composed of β pleated sheets and displayed in clusters of three subunits
      \textbf{Ligands}: TNF-family ligands are trimeric, β-sheet-rich proteins that are commonly synthesized as membrane-bound juxtacrine signals and may subsequently be released in soluble form.

      % Pathway binds to Death Domain (DD), TRADD, FADD, and RIP → canonical NF-κB pathway
      Receptor outcome depends on the family member and adaptor complex. TNFR1 can recruit TRADD and RIPK1 to activate canonical NF-κB and MAPK signaling, while alternative complexes involving FADD and caspase-8 can initiate apoptosis.

    \end{core}

    \begin{core}{Chemokines}

      <FIGURE>

      % Structure
      \textbf{Structure}: Chemokine receptors are seven-transmembrane G-protein-coupled receptors whose cytoplasmic regions interact with heterotrimeric G proteins.

      % 7-spanning α helices in membrane associated with G-protein on C-terminal
      Their seven membrane-spanning α helices transmit ligand-induced conformational changes to a G protein associated with the receptor's intracellular surface.

      % Very promiscuous receptor (cross reception)
      Chemokine recognition is frequently promiscuous: one receptor may bind several chemokines, and one chemokine may engage multiple receptors.

      % Ligands—α helices and β pleated sheets stabilized with CCCCs
      \textbf{Ligands}: Chemokines are small secreted proteins containing a flexible amino terminus, a three-stranded β sheet, and a carboxy-terminal α helix stabilized by conserved disulfide bonds.

      % G Protein Pathway—
      \textbf{G-protein pathway}: Ligand engagement activates an associated heterotrimeric G protein.

      % Ligand binding to receptor induces α subunit of G protein to switch out its GDP to GTP
      The activated receptor acts as a guanine-nucleotide exchange factor, causing the Gα subunit to release GDP and bind GTP.

      % α, β, γ G Protein subunits all dissociate and start their individual signaling cascades
      GTP-bound Gα separates functionally from the Gβγ dimer, and both components regulate downstream effectors controlling chemotaxis, adhesion, and cytoskeletal remodeling.

    \end{core}

  \subsection{Immune Signaling Types \& Themes}

    \begin{core}{Immune Signaling Types}

      <TABLE>

    \end{core}

    \begin{core}{Signaling Themes}

      <FIGURE>

      \begin{enumerate}

        \item \textbf{Phosphorylation} 

          % Adding phosphates to, and removing phosphates from molecules
          Reversible phosphorylation adds or removes phosphate groups from proteins or lipids, thereby altering their activity, localization, conformation, or molecular interactions.

          % Examples: BCRs, JAK-STAT, TCRs, G Proteins, Membrane inositols
          This regulatory principle is central to BCR and TCR signaling, JAK--STAT pathways, G-protein-coupled receptor pathways, and the metabolism of membrane phosphoinositides.

        \item \textbf{Phospholipid Hydrolysis}

          % Phospholipids are split
          Phospholipid hydrolysis generates membrane-associated and soluble second messengers from phospholipid substrates.

          % Examples:
          A prominent example is the cleavage of phosphatidylinositol 4,5-bisphosphate (PIP$_2$) by phospholipase C.

          % Phospholipid is split into
          This reaction produces two functionally distinct products:

          \begin{itemize}

            % DAG (non polar part)
            \item diacylglycerol (DAG), a nonpolar messenger that remains within the plasma membrane; and

            % IP3 (polar part) cascades
            \item inositol 1,4,5-trisphosphate (IP$_3$), a soluble messenger that diffuses through the cytosol.

          \end{itemize}

          % *IP3 makes the ER leaky to Ca2+, leading to the NFAT pathway through binding to calmodulin
          IP$_3$ binds its receptor on the RER and opens Ca$^{2+}$ channels. The resulting rise in cytosolic Ca$^{2+}$ activates calmodulin and calcineurin, which promote NFAT translocation into the nucleus.

        \item \textbf{Protein Hydrolysis}

          % Using hydrolysis of cytosolic components in situations where you don’t want to reverse the signal (once you’ve clipped a protein, you can’t reassemble it, you must synthesize a new one)
          Regulated proteolysis irreversibly cleaves a protein substrate and is therefore well suited to decisive signaling transitions. Restoration of the original state generally requires synthesis of a new intact protein.

          % Examples: Notch Signaling, Apoptosis Induction
          Examples include proteolytic release of the Notch intracellular domain and caspase-mediated cleavage of cellular substrates during apoptosis.

        \item \textbf{Cytoskeletal Changes}

          % In relation to the movement of a cell throughout the body
          Cytoskeletal remodeling enables immune cells to polarize, migrate, adhere, phagocytose targets, and organize receptor signaling interfaces.

          % Examples: Chemokine → 7-spanR → G Protein signal → Integrin conformational change → Talins → Actin microfilaments → Movement
          For example, chemokine engagement of a seven-transmembrane receptor activates heterotrimeric G proteins, induces inside-out integrin activation, recruits talin and other adaptor proteins, and couples adhesion receptors to actin remodeling and directed movement.

      \end{enumerate}

    \end{core}

  \subsection{Cytokine Interaction Types}

    \begin{core}{Cytokine Interaction Types}

      \begin{enumerate}

        \item \textbf{Pleiotropy}

          <FIGURE>

          % One cytokine may act on several different cell types
          \textbf{Pleiotropy} occurs when a single cytokine produces distinct effects in multiple target-cell populations.

          % IL-4 coordinates a Th2 response by acting on both B and mast cells
          For example, IL-4 coordinates a T$_{\text{H}}$2-associated response by acting on B cells, mast cells, macrophages, and helper T cells.

        \item \textbf{Redundancy}

          <FIGURE>

          % More than one cytokine will produce the same effect
          \textbf{Redundancy} occurs when different cytokines can elicit overlapping biological outcomes.

          % IL-2, 4, and 5 induces proliferation in B cells
          For example, IL-2, IL-4, and IL-5 can each support aspects of activated B-cell proliferation under appropriate conditions.

        \item \textbf{Synergy}

          <FIGURE>

          % Combinations of stimulus have effects neither has by itself
          \textbf{Synergy} occurs when two signals acting together generate a response greater than, or qualitatively different from, the response to either signal alone.

          % IL-4 and IL-5 together will cause B cells to class switch and make IgE immunoglobulins
          In a T$_{\text{H}}$2-associated response, IL-4 provides a principal signal for IgE class-switch recombination, while IL-5 can support the proliferation and differentiation of appropriately activated B cells.

        \item \textbf{Antagonism} 

          <FIGURE>

          % One cytokine block the effects of another
          \textbf{Antagonism} occurs when one cytokine inhibits or counterbalances the effects of another.

          % IFNγ (a Th1 promoter) blocks the ability of IL-4 (a Th2 promoter) to induce class switching to IgE
          For example, IFN-γ promotes T$_{\text{H}}$1-associated immunity and can oppose IL-4-driven T$_{\text{H}}$2 differentiation and IgE class switching.

        % \item \textbf{Cascade Inductionk}
        \item \textbf{Cascade Induction}

          <FIGURE>

          % One cytokine stimulates other cells to produce a different cytokine, which stimulates more cells
          \textbf{Cascade induction} occurs when one cytokine stimulates a target cell to produce additional cytokines, thereby propagating and amplifying the response.

          % Macrophage → IL-12 → Th1 → IFNγ → activated macrophage pathway of development
          For example, IL-12 produced by activated macrophages and dendritic cells promotes T$_{\text{H}}$1 differentiation; T$_{\text{H}}$1-derived IFN-γ then enhances classical macrophage activation, reinforcing the cell-mediated immune response.

      \end{enumerate}

    \end{core}
